\documentclass[11pt,oneside]{book}
\usepackage[english, italian]{babel} 
\usepackage{microtype}
\usepackage[utf8]{inputenc}
\usepackage{subfigure}
\usepackage{epsfig}
\usepackage{amsmath, amssymb}
\usepackage[linesnumbered,lined,commentsnumbered,italiano]{algorithm2e}
\usepackage{fancybox}
\usepackage{listings}
\usepackage{url}
\usepackage{graphicx}
\usepackage{todonotes}
\usepackage{moreverb}
\usepackage{lmodern}
\usepackage{multirow}
\usepackage[labelfont=it,textfont={it}]{caption}
\usepackage[final,bookmarks,breaklinks,colorlinks,allcolors=black]{hyperref}
\usepackage{color}
\usepackage[block=ragged,backend=biber,natbib=true]{biblatex} % Use the bibtex backend with the authoryear citation style (which resembles APA)
\addbibresource{bibliography.bib} % The filename of the bibliography
\usepackage[autostyle=true]{csquotes} % Required to generate language-dependent quotes in the bibliography
\usepackage{hyperref}

%%%%%%%%%%%%%%%%%%%%%%%%%%%%
\definecolor{mygreen}{rgb}{0,0.6,0}
\definecolor{mygray}{rgb}{0.5,0.5,0.5}
\definecolor{mymauve}{rgb}{0.58,0,0.82}
\lstset{ %
	backgroundcolor=\color{white},   % choose the background color; you must add \usepackage{color} or \usepackage{xcolor}; should come as last argument
	basicstyle=\footnotesize,        % the size of the fonts that are used for the code
	breakatwhitespace=false,         % sets if automatic breaks should only happen at whitespace
	breaklines=true,                 % sets automatic line breaking
	captionpos=b,                    % sets the caption-position to bottom
	commentstyle=\color{mygreen},    % comment style
	deletekeywords={...},            % if you want to delete keywords from the given language
	escapeinside={\%*}{*)},          % if you want to add LaTeX within your code
	extendedchars=true,              % lets you use non-ASCII characters; for 8-bits encodings only, does not work with UTF-8
	frame=single,
	keepspaces=true,                 % keeps spaces in text, useful for keeping indentation of code (possibly needs columns=flexible)
	keywordstyle=\color{blue},       % keyword style
	morekeywords={*,...},            % if you want to add more keywords to the set
	numbers=left,                    % where to put the line-numbers; possible values are (none, left, right)
	numberstyle=\tiny\color{mygray}, % the style that is used for the line-numbers
	rulecolor=\color{black},         % if not set, the frame-color may be changed on line-breaks within not-black text (e.g. comments (green here))
	showspaces=false,                % show spaces everywhere adding particular underscores; it overrides 'showstringspaces'
	showstringspaces=false,          % underline spaces within strings only
	showtabs=false,                  % show tabs within strings adding particular underscores
	stringstyle=\color{mymauve},     % string literal style
	tabsize=2,	                   % sets default tabsize to 2 spaces
	title=\lstname                   % show the filename of files included with \lstinputlisting; also try caption instead of title
}
%%%%%%%%%%%%%
\definecolor{gray}{rgb}{0.4,0.4,0.4}
\definecolor{darkblue}{rgb}{0.0,0.0,0.6}
\definecolor{cyan}{rgb}{0.0,0.6,0.6}

%%%%%%%%%%%%%%%%%%%%%%%%
\begin{document}
\selectlanguage{italian}

\begin{titlepage}
\begin{center}
\epsfig{file=Figure/logo_standard.jpg,width=2.5truecm}\\[0.2truecm]
{\Large Universit\`a degli Studi di Salerno}\\[0.2truecm]
{\large Dipartimento di Informatica}\\
\hrulefill
\vfill
{\Large GameUp }\\[0.2truecm]
\vfill\vfill
{\Huge Problem Statement}
\vfill\vfill


{\bf Studenti} \hfill {\bf Docente}\ \ \\
Francesco Foglia \hfill Prof. Andrea De Lucia\\

\vfill
\hrulefill 

Anno Accademico 2020-2021

\end{center}
\end{titlepage}

\pagenumbering{roman}
\chapter*{Cronologia Revisioni}
\begin{center}
	\begin{tabular}{||c c p{10cm}||} 
	\hline
	Data & Versione & Descrizione \\ [0.5ex] 
	\hline\hline
	24/01/2021 & 1.0 & Aggiunta problem statement mancante al set di deliverables \\ 
	\hline
   \end{tabular}
\end{center}

\tableofcontents
\pagestyle{plain}

%%%%%%%%%%%%%%%%%%%%%%%%%%%%
\chapter{Introduzione}
\setcounter{page}{1} 	% devono seguire solo il primo capitolo
\pagenumbering{arabic}	% devono seguire solo il primo capitolo
\input{01PS_introduzione.tex}

%%%%%%%%%%%%%%%%%%%%%%%%%%%%
\chapter{Requisiti funzionali e non funzionali}
I requisiti funzionali e non funzionali sono reperibili nel \emph{Requirements Analysis Document}.

%%%%%%%%%%%%%%%%%%%%%%%%%%%%
\chapter{Ambiente target}
Come descritto precedentemente, il target del sistema sarà composto dal mondo \emph{mobile}. Per poter offrire l'accesso alla piattaforma a quanti più tipi di dispositivi possibili, si adotterà una soluzione \emph{Web based}, così da rendere il sistema intrinsecamente portabile.

%%%%%%%%%%%%%%%%%%%%%%%%%%%%
\end{document}